% Options for packages loaded elsewhere
\PassOptionsToPackage{unicode}{hyperref}
\PassOptionsToPackage{hyphens}{url}
\PassOptionsToPackage{dvipsnames,svgnames,x11names}{xcolor}
%
\documentclass[
  11pt,
]{article}
\usepackage{amsmath,amssymb}
\usepackage{iftex}
\ifPDFTeX
  \usepackage[T1]{fontenc}
  \usepackage[utf8]{inputenc}
  \usepackage{textcomp} % provide euro and other symbols
\else % if luatex or xetex
  \usepackage{unicode-math} % this also loads fontspec
  \defaultfontfeatures{Scale=MatchLowercase}
  \defaultfontfeatures[\rmfamily]{Ligatures=TeX,Scale=1}
\fi
\usepackage{lmodern}
\ifPDFTeX\else
  % xetex/luatex font selection
\fi
% Use upquote if available, for straight quotes in verbatim environments
\IfFileExists{upquote.sty}{\usepackage{upquote}}{}
\IfFileExists{microtype.sty}{% use microtype if available
  \usepackage[]{microtype}
  \UseMicrotypeSet[protrusion]{basicmath} % disable protrusion for tt fonts
}{}
\makeatletter
\@ifundefined{KOMAClassName}{% if non-KOMA class
  \IfFileExists{parskip.sty}{%
    \usepackage{parskip}
  }{% else
    \setlength{\parindent}{0pt}
    \setlength{\parskip}{6pt plus 2pt minus 1pt}}
}{% if KOMA class
  \KOMAoptions{parskip=half}}
\makeatother
\usepackage{xcolor}
\usepackage[margin=1.05in]{geometry}
\setlength{\emergencystretch}{3em} % prevent overfull lines
\providecommand{\tightlist}{%
  \setlength{\itemsep}{0pt}\setlength{\parskip}{0pt}}
\setcounter{secnumdepth}{-\maxdimen} % remove section numbering
\usepackage{amsmath,amssymb,mathtools}
\usepackage{microtype}
\allowdisplaybreaks
\usepackage{needspace}
\ifLuaTeX
  \usepackage{selnolig}  % disable illegal ligatures
\fi
\usepackage{bookmark}
\IfFileExists{xurl.sty}{\usepackage{xurl}}{} % add URL line breaks if available
\urlstyle{same}
\hypersetup{
  pdftitle={Lattice Edgeworth Expansions for Total Luce Cycles},
  pdfauthor={Christopher D. Long},
  colorlinks=true,
  linkcolor={black},
  filecolor={Maroon},
  citecolor={Blue},
  urlcolor={blue},
  pdfcreator={LaTeX via pandoc}}

\title{Lattice Edgeworth Expansions for Total Luce Cycles}
\author{Christopher D. Long}
\date{Research note -- 11 October 2026}

\begin{document}
\maketitle

\section{Results and scope}\label{results-and-scope}

Let \(\sigma_n\) be the label-to-rank permutation of independent
exponential clocks with rates \(\theta_{n,i}\). Assume that, after a
common rescaling in each row,

\[
1\le\theta_{n,i}\le\kappa
\]

for a fixed \(\kappa<\infty\). Write

\[
K_n=\sum_{r=1}^n C_{n,r},\qquad
\mu_n=\mathbb E K_n,\qquad
L=\log n,\qquad
\ell=1+\log L,
\]

and define the second factorial cumulant

\[
c_n=\operatorname{Var}K_n-\mu_n.
\]

All asymptotic estimates below are uniform over rate arrays with a fixed
comparability bound. The definition of \(\ell\) is used only for
sufficiently large \(n\).

The earlier total-cycle note proved

\[
d_{\mathrm{TV}}(\mathcal L(K_n),\operatorname{Pois}(\mu_n))
=O_\kappa(\ell^2/L),\qquad
\mu_n=L+O_\kappa(\ell).
\tag{1}
\]

This error is already \(o(L^{-1/2})\), so it transfers the first lattice
Edgeworth expansion of the Poisson law. The fixed-order factorial
domination in Paper IV allows a stronger signed approximation and a
second-order expansion. No limiting profile is needed.

\Needspace{10\baselineskip}

\textbf{Theorem 1 (First lattice Edgeworth expansion).} Let

\[
\psi(u)=\frac12-\{u\},\qquad \{u\}=u-\lfloor u\rfloor.
\]

In particular \(\psi(k)=1/2\) at integers: this is the convention for a
right-continuous distribution function. Uniformly in \(x\in\mathbb R\),

\[
\begin{split}
\mathbb P\left\{\frac{K_n-\mu_n}{\sqrt L}\le x\right\}
={}&\Phi(x)+\frac{\varphi(x)}{\sqrt L}
\left[\frac{1-x^2}{6}+\psi(\mu_n+x\sqrt L)\right]
+O_\kappa(\ell^2/L).
\end{split}
\tag{2}
\]

The identical assertion holds with \(L\) replaced by \(\mu_n\) in the
denominator, prefactor, and lattice argument. The first smooth
correction is universal. The oscillating term is essential and has the
same order as the smooth correction.

\Needspace{10\baselineskip}

\textbf{Corollary 2 (Exact leading Kolmogorov constant).}

\[
\sqrt L\,d_{\mathrm K}\left(
\mathcal L\left(\frac{K_n-\mu_n}{\sqrt L}\right),N(0,1)
\right)
\longrightarrow\frac{2}{3\sqrt{2\pi}}.
\tag{3}
\]

The convergence is uniform in the comparable-rate class. The same limit
holds with scale \(\sqrt{\mu_n}\) or with the exact standard deviation.

To state the stronger expansion, let

\[
p_\mu(k)=e^{-\mu}\frac{\mu^k}{k!}\quad(k\ge0),\qquad
p_\mu(k)=0\quad(k<0),
\]

and use the right-shift difference \(\Delta p(k)=p(k-1)-p(k)\).

\Needspace{10\baselineskip}

\textbf{Theorem 3 (Second-order signed Poisson approximation).}

\[
c_n=O_\kappa(\ell^2),\qquad
\operatorname{Var}K_n=L+O_\kappa(\ell^2),
\tag{4}
\]

and

\[
\frac12\sum_{k\in\mathbb Z}
\left|\mathbb P(K_n=k)-p_{\mu_n}(k)
-\frac{c_n}{2}\Delta^2p_{\mu_n}(k)\right|
\le C_\kappa\frac{\ell^3}{L^{3/2}}.
\tag{5}
\]

The approximating sequence is a signed measure of total mass one;
positivity is not claimed or required.

\Needspace{10\baselineskip}

\textbf{Theorem 4 (Second-order continuity-corrected Edgeworth
expansion).} For \(k\in\mathbb Z\), set

\[
y=y_{n,k}=\frac{k+1/2-\mu_n}{\sqrt{\mu_n}},\qquad
Q(y)=\frac{-y^5+7y^3-3y}{72}.
\]

Uniformly in integer \(k\),

\[
\begin{split}
\mathbb P(K_n\le k)
={}&\Phi(y)+\frac{\varphi(y)}{6\sqrt{\mu_n}}(1-y^2)\\
&+\frac{\varphi(y)}{\mu_n}
\left[Q(y)-\frac{c_n y}{2}\right]
+O_\kappa\left(\frac{\ell^3}{L^{3/2}}\right).
\end{split}
\tag{6}
\]

Since \(\ell^3/L^{3/2}=o(L^{-1})\), this resolves the second-order term.
Its first non-Poisson coefficient is exactly
\(c_n=\operatorname{Var}K_n-\mathbb E K_n\). The estimate does not
assert that this coefficient has a limit.

\Needspace{10\baselineskip}

\textbf{Corollary 5 (Exact-variance standardization).} Put
\(s_n^2=\operatorname{Var}K_n\) and \(z=(k+1/2-\mu_n)/s_n\). Then,
uniformly in \(k\),

\[
\mathbb P(K_n\le k)
=\Phi(z)+\frac{\varphi(z)}{6s_n}(1-z^2)
+\frac{\varphi(z)}{s_n^2}Q(z)
+O_\kappa(\ell^3/L^{3/2}).
\tag{7}
\]

Thus, after using the exact mean and variance, the displayed smooth
coefficients are universal through second order. The mean and variance,
and consequently the lattice locations, still depend on the rate array.

These conclusions are deductions from the named switching inputs
recorded below. This note does not re-prove the whole endpoint
construction, certify every argument in Paper IV, or establish a
profile-dependent mod-Poisson limit.

\section{Imported inputs and the additional moment
information}\label{imported-inputs-and-the-additional-moment-information}

The repository is frozen at commit

\texttt{6178a99a8f2ee7b6416edd026e7780b9b8fa9fa0}.

The quantitative inputs are Paper IV's \texttt{prop:poly-block},
\texttt{lem:passive-Stein}, and \texttt{thm:poly-pp}. The additional
input beyond the preceding Berry--Esseen note is
\texttt{cor:factorial-domination}, proved by the rate-oriented
fixed-order multi-clock identity \texttt{thm:multiclock}. No arithmetic
order statistic or prime-power estimate is needed.

For any fixed accuracy exponent \(Q_0\), choose a constant
\(A=A(\kappa,Q_0)\) sufficiently large and put

\[
b=\lceil L^3\rceil,\qquad U=\lfloor n/L^A\rfloor,
\]

\[
J_n=\sum_{b<r\le U}C_{n,r},\qquad
D_n=K_n-J_n,\qquad
\lambda_n=\sum_{b<r\le U}\frac{n}{r(n-r)}.
\]

A good split producing a core \(r\)-cycle comes from a source cycle of
length greater than \(r+U\). It changes \(J_n\) by one and leaves
\(D_n\) unchanged: the source and residual cycles are both above \(U\),
and the short cycles are unaffected. The marked theorem therefore
supplies an independent pair

\[
Z_n\sim\operatorname{Pois}(\lambda_n),\qquad
\widetilde D_n\overset d=D_n
\]

such that

\[
d_{\mathrm{TV}}\bigl(\mathcal L(J_n,D_n),
\mathcal L(Z_n,\widetilde D_n)\bigr)
\le C_{\kappa,Q_0}L^{-Q_0}.
\tag{8}
\]

Write \(S_n=Z_n+\widetilde D_n\), \(\nu_n=\mathbb ED_n\), and
\(\mu_n^0=\lambda_n+\nu_n\). Separately, the exact mean-switching
identity and the polynomial intensity error give

\[
|\mathbb EJ_n-\lambda_n|\le C_{\kappa,Q_0}L^{-Q_0},
\qquad |\mu_n-\mu_n^0|\le C_{\kappa,Q_0}L^{-Q_0}.
\tag{9}
\]

This is not a deduction of moment convergence from total variation.

The fixed-order domination says that for every fixed integer \(q\ge1\),

\[
\mathbb E\mathcal C_n(r_1,\ldots,r_q)
\le\frac{C_{\kappa,q}}{r_1\cdots r_q},
\qquad r_1+\cdots+r_q\le n/2,
\tag{10}
\]

where \(\mathcal C_n\) counts ordered tuples of distinct cycles, also
when some lengths coincide. The rate-oriented rotation used in its proof
has only one positive clock coefficient in its logarithmic density
ratio. Splitting the sum by the position of that clock permits injective
reindexing by the corresponding label. This gives the selected
first-moment bound for every fixed \(q\) without requiring a product of
\(q\) untruncated swap factors to have a high moment.

\subsection{All fixed moments of the boundary
count}\label{all-fixed-moments-of-the-boundary-count}

\Needspace{10\baselineskip}

\textbf{Lemma 6.} For every fixed integer \(q\ge1\),

\[
\mathbb ED_n^q=O_{\kappa,q}(\ell^q),\qquad
\mathbb EK_n^q=O_{\kappa,q}(L^q).
\tag{11}
\]

\textbf{Proof.} Let \(W=\lfloor n/(2q)\rfloor\). For large \(n\),
\(b<U<W\). Let \(V_n\) count the omitted cycles with lengths at most
\(W\). Their length set has harmonic mass

\[
\sum_{r\le b}\frac1r+\sum_{U<r\le W}\frac1r=O_{\kappa,q}(\ell).
\]

For every \(j\le q\), a tuple of \(j\) such cycles has total length at
most \(n/2\), so (10) gives \(\mathbb E(V_n)_j=O_{\kappa,q}(\ell^j)\).
Expanding \(V_n^q\) in falling factorials proves
\(\mathbb EV_n^q=O(\ell^q)\). At most \(2q-1\) cycles are longer than
\(W\), because each has length at least \(W+1>n/(2q)\). Adding this
bounded count proves the first assertion. The identical argument with
all lengths at most \(W\), whose harmonic mass is \(O(L)\), proves the
second. \(\square\)

The same global moment bound holds for \(S_n\): the Poisson parameter is
\(O(L)\) and the mark has the moments just proved. Under a maximal
coupling of \(K_n\) and \(S_n\), (8) and Cauchy--Schwarz imply

\[
|\mathbb EK_n^j-\mathbb ES_n^j|
\le C_{\kappa,j}L^j L^{-Q_0/2}
\tag{12}
\]

for every fixed \(j\) for which the global \(2j\)th moments are invoked.
Thus the probability approximation can legitimately be upgraded to the
specific moment comparison needed here, after increasing \(Q_0\).

In particular, if

\[
b_n^*=\operatorname{Var}D_n-\nu_n,
\]

then (9), (11), and (12) show

\[
c_n=b_n^*+O_\kappa(L^{2-Q_0/2}).
\tag{13}
\]

Taking \(Q_0\ge12\) makes this error \(O(L^{-4})\). Since
\(|b_n^*|=O_\kappa(\ell^2)\), (4) follows. More generally the same
reasoning gives, for each fixed \(j\ge2\),

\[
\kappa_j^F(K_n)=O_{\kappa,j}(\ell^j),\qquad
\kappa_j(K_n)=\mu_n+O_{\kappa,j}(\ell^j),
\tag{14}
\]

where \(\kappa_j^F\) denotes a factorial cumulant. These are
polylogarithmic bounds, not \(O(1)\) cumulant corrections.

\section{Discrete Taylor expansion and Poisson
smoothing}\label{discrete-taylor-expansion-and-poisson-smoothing}

Let \(T\) be right shift and \(\Delta=T-I\). For every integer \(m\ge0\)
and fixed integer \(s\ge0\),

\[
T^mp=\sum_{j=0}^s\binom mj\Delta^jp
+\sum_{u=0}^{m-s-1}\binom{m-1-u}{s}T^u\Delta^{s+1}p,
\tag{15}
\]

with an empty remainder when \(m\le s\) and with \(\binom mj=0\) for
\(j>m\). The identity follows by repeated telescoping, or induction on
\(m\). Translation preserves the \(\ell^1\) norm, and the sum of the
remainder coefficients is \(\binom m{s+1}\). Hence

\[
\left\|T^mp-\sum_{j=0}^s\binom mj\Delta^jp\right\|_1
\le\binom m{s+1}\|\Delta^{s+1}p\|_1.
\tag{16}
\]

\subsection{Norms of Poisson
differences}\label{norms-of-poisson-differences}

For \(Z\sim\operatorname{Pois}(\lambda)\), define the Charlier
polynomial by

\[
C_j(k;\lambda)=\sum_{r=0}^j(-1)^{j-r}\binom jr(k)_r\lambda^{j-r}.
\]

Then

\[
\Delta^jp_\lambda(k)=\lambda^{-j}C_j(k;\lambda)p_\lambda(k).
\]

The generating identity

\[
\sum_{j\ge0}\frac{C_j(k;\lambda)}{j!}t^j=e^{-\lambda t}(1+t)^k
\]

and the Poisson generating function give

\[
\mathbb E\left[e^{-\lambda(s+t)}(1+s)^Z(1+t)^Z\right]
=e^{\lambda st}.
\]

Comparing coefficients proves \(\mathbb EC_j(Z;\lambda)^2=j!\lambda^j\).
Therefore

\[
\boxed{\|\Delta^jp_\lambda\|_1\le\sqrt{j!}\,\lambda^{-j/2}.}
\tag{17}
\]

This is an exact orthogonality calculation, not a central-limit
approximation.

\subsection{Retaining the second factorial
cumulant}\label{retaining-the-second-factorial-cumulant}

Let \(M\) be an independent nonnegative integer mark, put
\(\nu=\mathbb EM\), and let \(N\sim\operatorname{Pois}(\nu)\). Apply
(15) with \(s=2\), once with \(m=M\) and once with \(m=N\). The constant
and linear terms cancel. If
\(b=\mathbb E(M)_2-\nu^2=\operatorname{Var}M-\nu\), then

\[
p_\lambda*\mathcal L(M)
=p_{\lambda+\nu}+\frac b2\Delta^2p_\lambda+R,
\]

\[
\|R\|_1\le\frac{\mathbb E(M)_3+\nu^3}{6}\|\Delta^3p_\lambda\|_1.
\]

Moreover, telescoping the shift and averaging over \(N\) gives

\[
\|\Delta^2p_{\lambda+\nu}-\Delta^2p_\lambda\|_1
\le\nu\|\Delta^3p_\lambda\|_1.
\]

Thus

\[
\begin{split}
\frac12\left\|p_\lambda*\mathcal L(M)-p_{\lambda+\nu}
-\frac b2\Delta^2p_{\lambda+\nu}\right\|_1
\le\frac{\sqrt6}{2\lambda^{3/2}}
\left[\frac{\mathbb E(M)_3+\nu^3}{6}+\frac{|b|\nu}{2}\right].
\end{split}
\tag{18}
\]

Apply (18) with \(M=\widetilde D_n\). Equation (11) makes its right side
\(O_\kappa(\ell^3/L^{3/2})\). The coupling error (8) is smaller.
Replacing \(\mu_n^0\) by \(\mu_n\) uses (9) and the elementary Poisson
parameter coupling. Replacing \(b_n^*\) by \(c_n\) uses (13) and (17).
All these replacements are smaller than the displayed remainder when
\(Q_0\ge12\). This proves Theorem 3.

\section{The elementary Poisson lattice
expansion}\label{the-elementary-poisson-lattice-expansion}

The analytic step can be derived directly for a Poisson variable, so no
nonlattice Edgeworth theorem is being applied to an integer-valued law.

Let \(h=\mu^{-1/2}\) and \(x_j=(j-\mu)h\). With the probabilists'
Hermite polynomials \(H_j\), Fourier inversion of the centered Poisson
characteristic function gives

\[
\begin{split}
p_\mu(j)=h\varphi(x_j)\left[
1+\frac h6H_3(x_j)
+h^2\left\{\frac{H_4(x_j)}{24}+\frac{H_6(x_j)}{72}\right\}
\right]+r_\mu(j),
\end{split}
\tag{19}
\]

where

\[
\sum_{j\in\mathbb Z}|r_\mu(j)|\le C h^3.
\tag{20}
\]

Here is a remainder justification. On the inversion interval
\([-\pi/h,\pi/h]\), the characteristic function is

\[
\chi_\mu(t)=\exp\{\mu(e^{iht}-1-iht)\},\qquad
|\chi_\mu(t)|\le e^{-2t^2/\pi^2}.
\]

On \(|t|\le h^{-1/8}\), expansion of its logarithm and then its
exponential, with two derivatives, compares it with

\[
e^{-t^2/2}\left[1+\frac h6(it)^3
+h^2\left\{\frac{(it)^4}{24}+\frac{(it)^6}{72}\right\}\right]
\]

with an integrable error and integrable second derivative both
\(O(h^3)\). Indeed the discarded terms are bounded by \(Ch^3\) times a
fixed polynomial in \(|t|\) times \(e^{-ct^2}\). Outside the cutoff, the
same Gaussian bound makes the integrals of the characteristic function
and its first two derivatives smaller than every power of \(h\); the
cutoff can be made smooth. Fourier inversion and two integrations by
parts give \(|r_\mu(j)|\le Ch^4/(1+x_j^2)\), with exponentially small
endpoint contributions handled by the periodic uncentered integrand.
Summing the grid, whose spacing is \(h\), proves (20).

For a Schwartz function \(f\) and \(y=(k+1/2-\mu)h\), midpoint
Euler--Maclaurin summation gives, uniformly in \(k\),

\[
h\sum_{j\le k}f(x_j)
=\int_{-\infty}^{y}f(t)\,dt-\frac{h^2}{24}f'(y)+O_f(h^3).
\tag{21}
\]

Apply this to the terms in (19). The identity

\[
\int_{-\infty}^{y}H_j(t)\varphi(t)\,dt
=-H_{j-1}(y)\varphi(y)
\]

and \(\varphi'(y)=-y\varphi(y)\) yield

\[
\begin{split}
\mathbb P\{\operatorname{Pois}(\mu)\le k\}
={}&\Phi(y)+\frac{\varphi(y)}{6\sqrt\mu}(1-y^2)\\
&+\frac{\varphi(y)}\mu
\left[-\frac{H_3(y)}{24}-\frac{H_5(y)}{72}+\frac y{24}\right]
+O(\mu^{-3/2}).
\end{split}
\tag{22}
\]

The bracket is exactly \(Q(y)=(-y^5+7y^3-3y)/72\). The positive \(y/24\)
term is the second-order continuity correction. Omitting it is not valid
at order \(\mu^{-1}\).

The use of a lattice correction is classical; see Kolassa--McCullagh
{[}KM90{]}. Signed Poisson--Charlier approximation schemes are developed
more generally in {[}BKN14, CDMN20, MNV22{]}. The proof here uses the
explicit Poisson core and moment bounds rather than assuming convergence
of a mod-Poisson residue.

\section{Proofs of the displayed Edgeworth
formulas}\label{proofs-of-the-displayed-edgeworth-formulas}

Summing the second difference telescopes exactly:

\[
\sum_{j\le k}\Delta^2p_\mu(j)
=p_\mu(k)-p_\mu(k-1)
=-\frac{k-\mu}{\mu}p_\mu(k).
\tag{23}
\]

With \(y=(k+1/2-\mu)/\sqrt\mu\), (19) implies, uniformly in \(k\),

\[
p_\mu(k)-p_\mu(k-1)
=-\frac{y\varphi(y)}\mu+O(\mu^{-3/2}).
\tag{24}
\]

Combine (5), (22), and (24), using \(c_n=O_\kappa(\ell^2)\). This proves
(6), including its sign and uniform remainder.

For a real threshold \(a=\mu+x\sqrt\mu\), let \(k=\lfloor a\rfloor\).
Its midpoint coordinate is

\[
y=x+\frac{\psi(a)}{\sqrt\mu}.
\]

Expanding (22) to first order, and using (1), proves

\[
\mathbb P\left\{\frac{K_n-\mu_n}{\sqrt{\mu_n}}\le x\right\}
=\Phi(x)+\frac{\varphi(x)}{\sqrt{\mu_n}}
\left[\psi(\mu_n+x\sqrt{\mu_n})+\frac{1-x^2}{6}\right]
+O_\kappa(\ell^2/L).
\tag{25}
\]

For scale \(\sqrt L\), use the same calculation at \(x\sqrt{L/\mu_n}\).
The lattice argument is then exactly \(\mu_n+x\sqrt L\); it is not
approximated. Since \(\mu_n-L=O_\kappa(\ell)\), the normal scale change
costs \(O_\kappa(\ell/L)\), proving (2).

For reference, the full second-order sawtooth form at scale
\(\sqrt{\mu_n}\) is also available. Set

\[
a=\mu_n+x\sqrt{\mu_n},\quad f=\{a\},\quad s=1/2-f,
\quad B_2(f)=f^2-f+1/6.
\]

Then

\[
\begin{split}
\mathbb P\left\{\frac{K_n-\mu_n}{\sqrt{\mu_n}}\le x\right\}
={}&\Phi(x)+\frac{\varphi(x)}{\sqrt{\mu_n}}
\left[s-\frac{H_2(x)}6\right]\\
&+\frac{\varphi(x)}{\mu_n}
\left[-\frac{H_3(x)}{24}-\frac{H_5(x)}{72}
+\frac{sH_3(x)}6-\frac x2\{B_2(f)+c_n\}\right]\\
&+O_\kappa(\ell^3/L^{3/2}).
\end{split}
\tag{26}
\]

This formula includes the right side of every jump. Formula (6) is
usually simpler to use.

\subsection{Exact-variance scaling}\label{exact-variance-scaling}

Let \(s_n^2=\mu_n+c_n\). In (6), replace \(y\) by
\(z\sqrt{1+c_n/\mu_n}\). The first-order normal-scale correction is
\(c_n z\varphi(z)/(2\mu_n)\), which cancels the term
\(-c_n y\varphi(y)/(2\mu_n)\). All remaining errors are

\[
O_\kappa\left(\frac{1+|c_n|}{\mu_n^{3/2}}
+\frac{c_n^2}{\mu_n^2}\right)
=O_\kappa(\ell^3/L^{3/2}).
\]

Replacing the remaining prefactors by \(s_n^{-1}\) and \(s_n^{-2}\) has
the same smaller cost. This proves (7). The lattice thresholds are
unchanged throughout.

\subsection{Exact leading Kolmogorov
constant}\label{exact-leading-kolmogorov-constant}

The first-order correction in (2) satisfies

\[
\left|\varphi(x)\left[\psi(\mu_n+x\sqrt L)
+\frac{1-x^2}{6}\right]\right|
\le\varphi(x)\left[\frac12+\frac{|1-x^2|}{6}\right].
\]

The right side has maximum \(2/(3\sqrt{2\pi})\), attained at \(x=0\). To
verify this, by symmetry take \(x\ge0\). On \([0,1]\) its polynomial
factor is \((4-x^2)/6\) and its derivative, after removing the common
Gaussian factor, is \(x(x^2-6)/6\). On \([1,\infty)\) the corresponding
factor is \((x^2+2)/6\) and the derivative is \(-x^3/6\). Both ranges
are decreasing away from zero.

For the matching lower bound, choose an integer \(k_n\) nearest
\(\mu_n\) and put \(x_n=(k_n-\mu_n)/\sqrt L\). Then \(x_n\to0\) and the
lattice term at this exact threshold equals \(1/2\). The correction
tends to \((1/2+1/6)\varphi(0)=2/(3\sqrt{2\pi})\). The error in (2) is
\(o(L^{-1/2})\). This proves (3). The changes to the other stated scales
are smaller than first order by (1) and (4).

\section{An expansion at every fixed Poisson--Charlier
order}\label{an-expansion-at-every-fixed-poissoncharlier-order}

The preceding argument is not limited to the second difference. The
following formulation uses exact finite-\(n\) coefficients and does not
assume a limiting profile.

For a fixed integer \(m\ge1\), define

\[
a_{n,j}=[w^j]\left\{e^{-\mu_n w}\mathbb E(1+w)^{K_n}\right\},
\qquad 0\le j\le m.
\]

The coefficients are well-defined derivatives at zero of an entire
function. There is no assertion of uniform analytic convergence in a
fixed complex neighborhood. In particular,

\[
a_{n,0}=1,\quad a_{n,1}=0,\quad
a_{n,2}=\frac{\kappa_2^F(K_n)}2=\frac{c_n}{2},
\]

\[
a_{n,3}=\frac{\kappa_3^F(K_n)}6,
\qquad
a_{n,4}=\frac{\kappa_4^F(K_n)+3\kappa_2^F(K_n)^2}{24}.
\]

\Needspace{10\baselineskip}

\textbf{Theorem 7 (Fixed-order signed expansion).} For every fixed
\(m\),

\[
|a_{n,j}|\le C_{\kappa,m}\ell^j\quad(2\le j\le m)
\]

and

\[
\frac12\left\|\mathcal L(K_n)-
\sum_{j=0}^m a_{n,j}\Delta^jp_{\mu_n}\right\|_1
\le C_{\kappa,m}\frac{\ell^{m+1}}{L^{(m+1)/2}}.
\tag{27}
\]

The order \(m\) is fixed before choosing constants and cutoffs. No
uniformity for an increasing expansion order is asserted.

\textbf{Proof.} First work with \(S=Z_\lambda+M\) from (8), whose mean
is \(\mu^0=\lambda+\nu\), and put

\[
\alpha_j=[w^j]\{e^{-\nu w}\mathbb E(1+w)^M\}.
\]

Equation (11) implies \(|\alpha_j|=O_{\kappa,m}(\ell^j)\). Applying
(15)--(17) to the mark gives an expansion of \(p_\lambda*\mathcal L(M)\)
through order \(m\) with coefficients \(\mathbb E(M)_j/j!\) and
remainder \(O(\ell^{m+1}\lambda^{-(m+1)/2})\) in \(\ell^1\).

For the proposed approximation
\(\sum_{j\le m}\alpha_j\Delta^jp_{\lambda+\nu}\), use
\(p_{\lambda+\nu}=p_\lambda*\operatorname{Pois}(\nu)\) and expand each
shift of \(\Delta^jp_\lambda\) through order \(m-j\). The error is
bounded by

\[
\sum_{j=0}^m|\alpha_j|\frac{\nu^{m-j+1}}{(m-j+1)!}
\|\Delta^{m+1}p_\lambda\|_1
=O_{\kappa,m}(\ell^{m+1}\lambda^{-(m+1)/2}).
\]

The coefficients through degree \(m\) agree, because

\[
\left(\sum_{j\ge0}\alpha_j w^j\right)e^{\nu w}
=\mathbb E(1+w)^M
\]

as a formal Taylor identity to the required finite order.

Finally choose the accuracy exponent \(Q_0\) large in terms of \(m\).
Equations (9), (11), and (12) show that the coefficients of
\(e^{-\mu_n w}\mathbb E(1+w)^{K_n}\) and those of
\(e^{-\mu_n^0w}\mathbb E(1+w)^{S_n}\) differ, through degree \(m\), by
\(O_{\kappa,m}(L^{j-Q_0/2})\) in degree \(j\). Equation (17), the mean
parameter replacement, and (8) make their contributions smaller than the
stated remainder. This proves (27). \(\square\)

Combining (27) with higher versions of (19)--(22) yields Gaussian
lattice expansions at each fixed order, with exact finite-\(n\)
factorial-cumulant coefficients. This is different from proving that
those coefficients converge to profile-dependent constants.

\section{Profile-dependent terms: deductions and unresolved
limits}\label{profile-dependent-terms-deductions-and-unresolved-limits}

There are three distinct questions: the location parameter, the first
non-Poisson shape coefficient, and a full limiting analytic residue. The
transfer argument identifies where each would enter, but does not by
itself identify its limit.

\subsection{A profile mean correction}\label{a-profile-mean-correction}

Let \(\delta_n=\mu_n-L=O_\kappa(\ell)\). Applying (2) at a shifted real
threshold gives the unconditional identity

\[
\begin{split}
\mathbb P\left\{\frac{K_n-L}{\sqrt L}\le x\right\}
={}&\Phi(x)+\frac{\varphi(x)}{\sqrt L}
\left[\psi(L+x\sqrt L)+\frac{1-x^2}{6}-\delta_n\right]
+O_\kappa(\ell^2/L).
\end{split}
\tag{28}
\]

For a sampled profile \(\theta_{n,i}=f(i/n)\) with \(f\in C([0,1])\)
strictly positive, an additional theorem \(\delta_n\to\kappa_f\) would
therefore yield

\[
\begin{split}
\mathbb P\left\{\frac{K_n-L}{\sqrt L}\le x\right\}
={}&\Phi(x)+\frac{\varphi(x)}{\sqrt L}
\left[\psi(L+x\sqrt L)+\frac{1-x^2}{6}-\kappa_f\right]
+o(L^{-1/2}).
\end{split}
\tag{29}
\]

The historical candidate proposed

\[
\kappa_f=\gamma_E+\sum_{r\ge1}\left(\lambda_r(f)-\frac1r\right),
\qquad \lambda_r(f)=\frac{\operatorname{Tr}(\mathcal K_f^r)}r.
\tag{30}
\]

This note does not reinstate the invalidated proof of that expectation
expansion. Formula (29) is conditional on the mean limit, and (30) is a
proposed identification, not a consequence of (27).

\subsection{A profile variance-minus-mean
correction}\label{a-profile-variance-minus-mean-correction}

If \(c_n\to c_f\), then (6) holds with \(c_f\) in place of \(c_n\) and
an \(o(L^{-1})\) remainder. This gives the profile shape contribution

\[
-\frac{c_f}{2\mu_n}y\varphi(y).
\tag{31}
\]

The proven general bound is only \(c_n=O_\kappa(\ell^2)\). It neither
proves convergence nor shows whether \(c_f\) is genuinely profile
dependent. The exact variance standardization in (7) absorbs this term.

Small changes inside the fractional-part function cannot be ignored
uniformly at its discontinuities. Approximate deterministic centers must
therefore be handled by their actual threshold, as in (28)--(29), not by
silently substituting an approximate mean into the lattice phase.

\subsection{A stronger, explicitly conjectural mod-Poisson
candidate}\label{a-stronger-explicitly-conjectural-mod-poisson-candidate}

For uniform permutations,

\[
\mathbb E z^{K_n}=\frac{\Gamma(n+z)}{\Gamma(z)\Gamma(n+1)},
\]

so the normalized generating function tends to \(1/\Gamma(z)\). This
classical example is discussed by Kowalski--Nikeghbali {[}KN10{]}. It
yields \(\mu_n=H_n\) and \(c_n=-H_n^{(2)}\to-\zeta(2)\).

Let \(A_f=\sum_r(\lambda_r(f)-1/r)\). A possible extension, suggested by
the short-cycle mean corrections and the uniform residue, is

\[
e^{-(z-1)L}\mathbb E z^{K_n}
\longrightarrow \frac{e^{A_f(z-1)}}{\Gamma(z)}
\tag{32}
\]

locally uniformly in a complex neighborhood of \(z=1\). This is a
\textbf{conjectural strengthening}, not an established Luce theorem. It
is not justified simply by fixed-dimensional Poisson convergence and
macroscopic Poisson--Dirichlet convergence.

If (32) were proved with the stated analytic convergence,
differentiating its logarithm would give

\[
\kappa_f=\gamma_E+A_f,
\qquad
\kappa_j^F(K_n)\longrightarrow
(-1)^{j+1}(j-1)!\zeta(j),\quad j\ge2.
\tag{33}
\]

Thus under this particular candidate, higher factorial-cumulant limits
would be universal rather than profile dependent. It predicts
\(c_f=-\zeta(2)\) for every regular profile. The finite-order expansion
theorem (27) does not prove (32); its constants depend on the order and
its coefficients need not converge.

\section{Source records, references, and finite
checks}\label{source-records-references-and-finite-checks}

\subsection{Internal dependencies}\label{internal-dependencies}

All repository line numbers refer to the frozen commit stated above.

\begin{itemize}
\tightlist
\item
  Paper IV, \texttt{luce-edgeworth.tex}: rate-oriented fixed-order
  switching and factorial domination, source lines 929--1112; polynomial
  one-block switching and passive-mark theorem, source lines 1840--2010.
  These modules were read directly for this note.
\item
  The preceding conversation artifact,
  \texttt{luce-total-cycle-berry-esseen.md}: the polynomial passive-mark
  coupling, exact-mean bookkeeping, and the second-moment Poisson
  smoothing theorem. The new note adds fixed-order moment transfer and
  retains the second and higher Poisson differences.
\item
  Historical provenance candidate
  \texttt{mesoscopic-total-cycle-resolution-candidate.tex}, source lines
  125--215: the proposed expectation constant in (30). This source is
  used only to identify the conjectural target and is not a proof
  dependency.
\end{itemize}

The prior functional audit and the preceding Berry--Esseen note pin the
same repository commit. No repository file is modified by this note.

\subsection{Primary literature}\label{primary-literature}

{[}KM90{]} J. E. Kolassa and P. McCullagh, \emph{Edgeworth Series for
Lattice Distributions}, Annals of Statistics \textbf{18} (1990),
981--985. DOI:
\href{https://doi.org/10.1214/aos/1176347637}{10.1214/aos/1176347637}.

{[}BKN14{]} A. D. Barbour, E. Kowalski, and A. Nikeghbali,
\emph{Mod-discrete expansions}, Probability Theory and Related Fields
\textbf{158} (2014), 859--893. DOI:
\href{https://doi.org/10.1007/s00440-013-0498-8}{10.1007/s00440-013-0498-8}.

{[}KN10{]} E. Kowalski and A. Nikeghbali, \emph{Mod-Poisson convergence
in probability and number theory}, International Mathematics Research
Notices \textbf{2010}, 3549--3587. DOI:
\href{https://doi.org/10.1093/imrn/rnq019}{10.1093/imrn/rnq019}. See in
particular the uniform-permutation formula (4.9).

{[}CDMN20{]} R. Chhaibi, F. Delbaen, P.-L. Méliot, and A. Nikeghbali,
\emph{Mod-\(\phi\) convergence: approximation of discrete measures and
harmonic analysis on the torus}, Annales de l'Institut Fourier
\textbf{70} (2020), 1115--1197. DOI:
\href{https://doi.org/10.5802/aif.3332}{10.5802/aif.3332}.

{[}MNV22{]} P.-L. Méliot, A. Nikeghbali, and G. Visentin,
\emph{Mod-Poisson approximation schemes and higher-order Chen--Stein
inequalities}, preprint
\href{https://arxiv.org/abs/2210.13818}{arXiv:2210.13818}.

\subsection{Verification performed}\label{verification-performed}

The accompanying \texttt{check\_edgeworth.py} checks 125 exact discrete
Taylor identities, six symbolic Charlier norm identities, the midpoint
polynomial and Kolmogorov envelope algebra, the Poisson CDF formulas for
eight parameters and five lattice phases, 30 independent
Poisson-plus-mark examples against the proved signed smoothing bound,
and exact uniform-permutation cycle distributions for
\(n=100,1000,10000,100000\).

All checks passed. Numerical computations test signs, phases,
coefficients, and finite consistency; they do not prove asymptotics or
independently verify the upstream Luce endpoint theorem. The JSON file
records the actual errors and test parameters.

\end{document}
